\chapter{Rules of calculus}
\label{chap:rules}

\begin{theorem}[Bilinear maps]
  \label{thm:bilinear}
  \lean{LCAnalysis.hasDirDerivAt_bilin, LCAnalysis.contDiffMBOn_bilin}
  \leanok
  \uses{def:ck, thm:partials}
  Let $\beta \colon A \times B \to F$ be continuous and bilinear. Then $\beta$ is $\Ck{k}$ for
  every $k$, with $\dd \beta\bigl((a, b), (h, w)\bigr) = \beta(a, w) + \beta(h, b)$.
\end{theorem}

\begin{proof}
  \leanok
  The difference quotient of $\beta$ at $(a,b)$ in the direction $(h,w)$ is
  $\beta(a, w) + \beta(h, b) + t\,\beta(h, w)$, which gives the formula. The differential is a
  sum of continuous bilinear maps composed with continuous linear maps, and induction on $k$
  completes the proof. This is G--N~1.2.3 and 1.3.5.
\end{proof}

\begin{corollary}[Product rule]
  \label{cor:product}
  \lean{LCAnalysis.ContDiffMBOn.bilin, LCAnalysis.ContDiffMBOn.dMB_bilin}
  \leanok
  \uses{thm:bilinear, thm:chain-rule, lem:prod}
  If $f \colon U \to A$ and $g \colon U \to B$ are $\Ck{k}$ and $\beta$ is continuous and
  bilinear, then $x \mapsto \beta(f(x), g(x))$ is $\Ck{k}$ and
  \[
    \dd\bigl(\beta(f, g)\bigr)(x)(h) = \beta\bigl(f(x), \dd g(x)(h)\bigr)
      + \beta\bigl(\dd f(x)(h), g(x)\bigr).
  \]
\end{corollary}

\begin{proof}
  \leanok
  Chain rule applied to $\beta \circ (f, g)$.
\end{proof}

\begin{lemma}[Derivatives along a line]
  \label{lem:line}
  \lean{LCAnalysis.ContDiffMBOn.hasDerivAt_iteratedDMB_line, LCAnalysis.ContDiffMBOn.continuousOn_iteratedDMB_line}
  \leanok
  \uses{thm:iterated-regular}
  Let $f$ be $\Ck{k}$ on $U$, $j < k$ and $\psi$ a continuous linear functional on $F$. Then the
  real function $s \mapsto \psi\bigl(\dd^{(j)} f(x + sh; h, \dots, h)\bigr)$ has derivative
  $\psi\bigl(\dd^{(j+1)} f(x + sh; h, \dots, h)\bigr)$ wherever $x + sh \in U$.
\end{lemma}

\begin{proof}
  \leanok
  By the chain rule along the line $s \mapsto (x + sh, h, \dots, h)$, whose derivative is
  $(h, 0, \dots, 0)$.
\end{proof}

\begin{theorem}[Taylor's formula, weak form]
  \label{thm:taylor}
  \lean{LCAnalysis.ContDiffMBOn.taylor}
  \leanok
  \uses{lem:line}
  Let $f$ be $\Ck{k+1}$ on a set containing the segment $[x, x+h]$. For every continuous linear
  functional $\psi$ on $F$,
  \[
    \psi\bigl(f(x + h)\bigr) = \sum_{j=0}^{k} \frac{1}{j!}\,
      \psi\bigl(\dd^{(j)} f(x; h, \dots, h)\bigr)
      + \int_0^1 \frac{(1-t)^k}{k!}\,
      \psi\bigl(\dd^{(k+1)} f(x + th; h, \dots, h)\bigr)\, dt .
  \]
\end{theorem}

\begin{proof}
  \leanok
  Apply Taylor's formula with integral remainder to the real function
  $s \mapsto \psi(f(x + sh))$, whose derivatives are given by Lemma~\ref{lem:line}. Compare
  G--N~1.6.1 and 1.6.14, where the formula is stated for the integrals in $F$; the weak form
  needs no completeness of $F$.
\end{proof}

\begin{theorem}[Second-order chain rule]
  \label{thm:chain-second}
  \lean{LCAnalysis.dMB_dMB_comp}
  \leanok
  \uses{thm:chain-rule}
  Let $\tau \colon U \to V$ and $g \colon V \to F$ be $\Ck{2}$, with $U \subseteq E$ and
  $V \subseteq E'$ open. For $y \in U$ and $a, b \in E$,
  \[
    \dd^{(2)} (g \circ \tau)(y; a, b)
      = \dd^{(2)} g\bigl(\tau y; \dd\tau(y)(a), \dd\tau(y)(b)\bigr)
        + \dd g(\tau y)\bigl(\dd^{(2)} \tau(y; a, b)\bigr).
  \]
\end{theorem}

\begin{proof}
  \leanok
  On $U \times E$, the first-order chain rule gives
  $\dd(g \circ \tau) = \dd g \circ \Phi$ with $\Phi(y', b') = (\tau y', \dd\tau(y')(b'))$.
  Differentiate this once more with the first-order chain rule. The derivative of $\Phi$ at
  $(y, b)$ in the direction $(a, 0)$ is $(\dd\tau(y)(a), \dd^{(2)}\tau(y; a, b))$. Split it as
  $(\dd\tau(y)(a), 0) + (0, \dd^{(2)}\tau(y; a, b))$: the first term gives
  $\dd^{(2)} g$, and the second gives $\dd g$, because $\dd g(\tau y)$ is linear.
\end{proof}
